\documentclass[11pt]{article}
\usepackage[margin=1in]{geometry}
\usepackage{enumitem}
% =============================================================================
% Preambles/header.tex — shared LaTeX/Beamer preamble for this project
%
% Use from a Beamer lecture by prepending:
%     \input{header}
% and compiling with TEXINPUTS=../Preambles:$TEXINPUTS (the /compile-latex
% skill does this automatically).
%
% PALETTE CONTRACT. Color names defined here MUST match the names in
% Quarto/theme-template.scss so Beamer and Quarto renderings use the same
% palette. The scripts/check-palette-sync.sh script enforces this.
%
% To customize: edit the HEX values below AND the matching $variable in
% Quarto/theme-template.scss, then run scripts/check-palette-sync.sh.
% =============================================================================

% -----------------------------------------------------------------------------
% Engine detection + encoding
%
% Our /compile-latex skill uses XeLaTeX, where inputenc is unnecessary and
% can produce encoding warnings. Only load inputenc under pdfTeX.
% -----------------------------------------------------------------------------
\usepackage{iftex}
\ifPDFTeX
  \usepackage[utf8]{inputenc}
\fi

\usepackage{amsmath, amssymb, amsthm}
\usepackage{booktabs}
\usepackage{graphicx}

% -----------------------------------------------------------------------------
% PALETTE — mirrors Quarto/theme-template.scss variable names.
% Load xcolor and define colors BEFORE anything that references them
% (notably hyperref, hypersetup, and the Beamer color assignments).
% -----------------------------------------------------------------------------
\usepackage{xcolor}

\definecolor{primary-blue}{HTML}{012169}      % headings, accents
\definecolor{primary-gold}{HTML}{B9975B}      % emphasis, borders
\definecolor{highlight-yellow}{HTML}{F2A900}  % markers, alerts
\definecolor{light-bg}{HTML}{E8EDF5}          % title / section backgrounds
\definecolor{jet}{HTML}{1A1A1A}               % body text

% Semantic colors (used in diagrams and annotations)
\definecolor{positive}{HTML}{15803D}          % good / observed / identified
\definecolor{negative}{HTML}{B91C1C}          % bad / problematic / bias
\definecolor{neutral}{HTML}{525252}           % reference / context

% Highlight accents (match .hi-* classes in the SCSS)
\definecolor{hi-slate}{HTML}{314F4F}
\definecolor{hi-green}{HTML}{2E8B57}
\definecolor{hi-red}{HTML}{C41E3A}

% Semantic aliases so skill templates and snippets can write
% \textcolor{accent}{...} without hard-coding "primary-blue".
\colorlet{accent}{primary-blue}
\colorlet{accent2}{primary-gold}

% -----------------------------------------------------------------------------
% Hyperref — loaded AFTER xcolor + palette so link colors resolve.
% -----------------------------------------------------------------------------
\usepackage{hyperref}
\hypersetup{colorlinks=true, linkcolor=primary-blue, urlcolor=primary-blue,
            citecolor=primary-blue, pdfborder={0 0 0}}

% -----------------------------------------------------------------------------
% Course code — set once per deck (after \input{header}, before \begin{document})
% via \coursecode{CS401}. Shown in the footer on every slide so a deck's course
% is identifiable without opening the title page. Empty by default.
% -----------------------------------------------------------------------------
\makeatletter
\newcommand{\coursecode}[1]{\gdef\@coursecodetext{#1}}
\gdef\@coursecodetext{}
\makeatother

% -----------------------------------------------------------------------------
% Beamer theme assignments (only apply under Beamer).
%
% We detect Beamer via \@ifclassloaded{beamer}, which is the documented,
% non-internal way. This must be wrapped in \makeatletter/\makeatother
% because \@ifclassloaded contains an @-catcode character.
% -----------------------------------------------------------------------------
\makeatletter
\@ifclassloaded{beamer}{%
  \usefonttheme{professionalfonts}
  \setbeamercolor{structure}{fg=primary-blue}
  \setbeamercolor{frametitle}{fg=primary-blue}
  \setbeamercolor{title}{fg=primary-blue}
  \setbeamercolor{subtitle}{fg=primary-gold}
  \setbeamercolor{author}{fg=primary-blue}
  \setbeamercolor{institute}{fg=neutral}
  \setbeamercolor{date}{fg=primary-gold}
  \setbeamercolor{itemize item}{fg=primary-blue}
  \setbeamercolor{itemize subitem}{fg=primary-gold}
  \setbeamercolor{alerted text}{fg=primary-gold}
  \setbeamercolor{block title}{fg=white, bg=primary-blue}
  \setbeamercolor{block body}{bg=light-bg}
  % Semantic block variants: block=definition/concept (blue), exampleblock=
  % worked example (green/positive), alertblock=key takeaway or caution (gold).
  % Keep to <=2 colored boxes per slide across ALL three variants (INV-7).
  \setbeamercolor{block title example}{fg=white, bg=positive}
  \setbeamercolor{block body example}{bg=light-bg}
  \setbeamercolor{block title alerted}{fg=white, bg=primary-gold}
  \setbeamercolor{block body alerted}{bg=light-bg}

  % Minimal footer: course code (left) + page X / N (right), no navigation clutter
  \setbeamertemplate{navigation symbols}{}
  \setbeamertemplate{footline}{%
    \color{neutral}\footnotesize\hspace{0.7em}\@coursecodetext\hfill
    \insertframenumber/\inserttotalframenumber\hspace{0.7em}\vspace{0.3em}}
}{}
\makeatother

% -----------------------------------------------------------------------------
% TikZ setup — libraries the snippet gallery uses
% -----------------------------------------------------------------------------
\usepackage{tikz}
\usetikzlibrary{arrows.meta, positioning, calc, decorations.pathreplacing,
                fit, shapes.geometric, backgrounds}

% Shared TikZ styles — snippets and hand-written diagrams can pick these up
% via `\begin{tikzpicture}[every node/.style={...}]` or by naming specific
% styles (e.g., `\node[dag-node] (X) at (0,0) {X};`).
\tikzset{
  % Node styles for causal DAGs and similar
  dag-node/.style = {
    draw=primary-blue, fill=light-bg, rounded corners=2pt,
    minimum width=1.2cm, minimum height=0.9cm, align=center, font=\small
  },
  decision-node/.style = {
    draw=primary-gold, fill=white, diamond, aspect=1.8,
    minimum width=1.8cm, minimum height=1.2cm, align=center, font=\small,
    inner sep=1pt
  },
  % Edge styles — observed vs counterfactual
  observed-edge/.style = {-{Latex[length=2.5mm]}, thick, draw=jet},
  counterfactual-edge/.style = {-{Latex[length=2.5mm]}, thick, draw=neutral, dashed},
  confound-edge/.style = {-{Latex[length=2.5mm]}, thick, draw=negative, dashed},
  % Data-point styles for plots
  observed-dot/.style = {fill=primary-blue, draw=primary-blue, circle, inner sep=1.2pt},
  counterfactual-dot/.style = {fill=white, draw=neutral, circle, inner sep=1.2pt}
}

% -----------------------------------------------------------------------------
% Convenience macros
% -----------------------------------------------------------------------------
\newcommand{\muted}[1]{\textcolor{neutral}{#1}}
\newcommand{\key}[1]{\textcolor{primary-gold}{\textbf{#1}}}
\newcommand{\good}[1]{\textcolor{positive}{#1}}
\newcommand{\bad}[1]{\textcolor{negative}{#1}}

% Standout frame macro: full-bleed dark background for section transitions.
% Beamer-only (uses \begin{frame}/\setbeamercolor/\usebeamerfont) -- do not
% call from an article-class file (e.g. Notes/<CODE>/*-notes.tex). \muted,
% \key, \good, \bad, and \coursecode above are plain xcolor/gdef and are
% safe to use from either class; verified when Notes/article-class support
% was added.
% Expand: \transitionslide{Your transition title}
\newcommand{\transitionslide}[1]{%
  \begingroup
  \setbeamercolor{background canvas}{bg=primary-blue}
  \begin{frame}[plain]
    \centering
    \vfill
    {\usebeamerfont{title}\color{white}\Large #1\par}
    \vfill
  \end{frame}
  \endgroup
}

% Section divider frame (Beamer-only), an alternative to \transitionslide with a
% darker two-line style: near-black (jet) background, a small white label line
% above a larger gold title, and a thin gold rule along the bottom edge.
% Expand: \sectiondivider{Section 2}{Pointers}
\newcommand{\sectiondivider}[2]{%
  \begingroup
  \setbeamercolor{background canvas}{bg=jet}
  \begin{frame}[plain]
    \vspace*{3.0cm}
    {\color{white}\ttfamily\large #1\par}
    \vspace{0.5cm}
    {\color{primary-gold}\LARGE #2\par}
    \begin{tikzpicture}[remember picture, overlay]
      \draw[primary-gold, line width=2.5pt]
        ($(current page.south west)+(0,0.1)$) -- ($(current page.south east)+(0,0.1)$);
    \end{tikzpicture}
  \end{frame}
  \endgroup
}

% -----------------------------------------------------------------------------
% Channel watermark — "Concepts That Click" (YouTube).
%
% Article-class files (CompetitiveExam Q/A sets, Notes, Assignments) get a
% light diagonal draft watermark on every page plus a centered footer naming
% the channel. Beamer decks get a subtle TikZ background watermark instead
% (the footline already carries the course code). Centralized here so every
% MCQ PDF is branded without per-file edits.
% -----------------------------------------------------------------------------
\makeatletter
\@ifclassloaded{article}{%
  \usepackage{draftwatermark}
  \SetWatermarkText{Concepts That Click}
  \SetWatermarkScale{1}
  \SetWatermarkFontSize{2cm}
  \SetWatermarkLightness{0.86}
  \SetWatermarkAngle{45}
  \usepackage{fancyhdr}
  \pagestyle{fancy}
  \fancyhf{}
  \fancyfoot[C]{\footnotesize\textcolor{neutral}{Concepts That Click~|~YouTube: \href{https://www.youtube.com/@concepts.thatclick}{@concepts.thatclick}}}
  \renewcommand{\headrulewidth}{0pt}
  \renewcommand{\footrulewidth}{0pt}
  \fancypagestyle{plain}{%
    \fancyhf{}%
    \fancyfoot[C]{\footnotesize\textcolor{neutral}{Concepts That Click~|~YouTube: \href{https://www.youtube.com/@concepts.thatclick}{@concepts.thatclick}}}%
    \renewcommand{\headrulewidth}{0pt}%
    \renewcommand{\footrulewidth}{0pt}%
  }
}{}
\@ifclassloaded{beamer}{%
  \setbeamertemplate{background}{%
    \begin{tikzpicture}[remember picture, overlay]
      \node[rotate=45, opacity=0.055, align=center]
        at (current page.center)
        {\Huge\textcolor{neutral}{Concepts That Click}};
    \end{tikzpicture}%
  }
}{}
\makeatother 
\coursecode{JKSSB}
\title{Lesson 19 Practice: Answers and Rationales}
\author{JKSSB Computer Awareness}
\date{\today}
\begin{document}
\maketitle
\noindent\textbf{Provenance:} Q1--Q3 are real JKSSB PYQs from the Website
Operator 2026 paper (\textit{[PYQ WO 2026 Q61--Q63, key:provisional]}); the
remaining 17 are original JKSSB-pattern questions (\textit{[Original]}).
\begin{enumerate}[label=\textbf{Q\arabic*.}, leftmargin=*]
\item \textbf{B.} The operating system provides a user-friendly interface
  layer between the user and the computer hardware. A calls it hardware when it
  is software; C calls it an application when it is system software; D calls it
  storage, which it is not.
\item \textbf{C.} While running, the OS is resident in primary storage (main
  memory). Secondary storage (A) and offline tape (B) are non-volatile stores;
  the display adapter (D) is an output device, not where the OS runs.
\item \textbf{D.} Multitasking means running more than one program at a time.
  Connecting several printers (A), storing several files (B), and using several
  monitors (C) are unrelated to multitasking.
\item \textbf{C.} OS stands for Operating System. A, B, and D are not the
  correct expansions of the abbreviation.
\item \textbf{A.} The kernel is the core part of the OS that interacts directly
  with the hardware. The GUI (B) is just the interface; the bootstrap (C) is a
  startup loader; a compiler (D) is a language tool, not the OS core.
\item \textbf{D.} Designing web pages is an application task, not an OS
  function. Process (A), memory (B), and file (C) management are all core OS
  functions.
\item \textbf{D.} GUI stands for Graphical User Interface. A, B, and C are not
  the correct full forms.
\item \textbf{B.} A CLI (Command Line Interface) takes typed text commands. A
  GUI (A) uses icons and clicks; the Ribbon (C) is a GUI toolbar; the Desktop
  (D) is a GUI workspace.
\item \textbf{A.} Starting the computer and loading the OS into main memory is
  called booting. Formatting (B) prepares storage; compiling (C) translates
  source code; defragmenting (D) reorganises disk data.
\item \textbf{C.} BIOS/UEFI initialises the hardware and then loads the OS
  during startup. A device driver (A) controls one device; the kernel (B) is
  the OS core loaded by the boot process; a compiler (D) translates programs.
\item \textbf{A.} A program in execution is a process. A folder (B), a shortcut
  (C), and a file (D) are stored items, not running programs.
\item \textbf{B.} Memory management decides how much main memory each program
  gets. Process management (A) schedules the CPU; file management (C) handles
  files; device management (D) controls input/output devices.
\item \textbf{C.} File management creates, reads, writes, and deletes files and
  folders. Scheduling (A), memory allocation (B), and print handling (D) belong
  to process, memory, and device management respectively.
\item \textbf{C.} A device driver lets the OS communicate with a specific
  hardware device. A compiler (A) translates code; a browser (B) is an
  internet application; a spreadsheet (D) is an application program.
\item \textbf{A.} The kernel is the core of the operating system, so this
  pairing is correct. A device driver (B) is not a text editor; the bootstrap
  (C) is not storage; a GUI (D) is not a command-line-only interface.
\item \textbf{B.} Time-sharing lets many users use a single computer at the
  same time. A computer cannot work without an OS (A), files need a file
  system (C), and hardware needs power (D).
\item \textbf{C.} Android is the most common example of an open-source
  operating system and is built on the Linux kernel. Windows 11 (A) and macOS
  (B) are proprietary; MS Word (D) is an application, not an OS.
\item \textbf{A.} Statements I and II are correct: the OS provides a
  user-friendly interface, and the kernel is its core. III is false because the
  OS is system software, not an application. B, C, and D all include the false
  statement III.
\item \textbf{C.} Logging in with a username and password is authentication, a
  part of security management. Compression (A), defragmentation (B), and
  spooling (D) are unrelated to login.
\item \textbf{B.} Updates fix security holes and add improvements. They cannot
  enlarge the physical disk (A), turn a GUI into hardware (C), or remove the
  need for a kernel (D).
\end{enumerate}
\end{document}
